\chapter{State of the Art and Theoretical Background}

\section{Introduction}
The design of an intelligent, safe, and conversational orchestration agent for enterprise firewalls requires a rigorous mastery of two intertwined disciplines: modern network administration automation and the controlled integration of artificial intelligence in critical systems. The primary objective of this chapter is to establish a structured theoretical foundation and conduct a comparative analysis of available technologies, culminating in a formally justified selection of the technical stack adopted in this project.

First, we introduce the fundamental theoretical concepts that underpin the project, covering the historical evolution of network administration paradigms, the specific engineering challenges of stateful firewall automation, and the role of key AI constructs — LLMs, RAG, and HITL safeguards. We then contrast the proposed hybrid orchestration model against existing paradigms to motivate the architectural direction. We follow with a detailed presentation of the generic architecture and functioning of a Hybrid AI Orchestration Agent, establishing a clear conceptual model that bridges theory to implementation. Finally, we conduct a structured comparative analysis of available technical solutions, formally justifying the technology stack selected for this project.

%% -----------------------------------------------------------
\section{Fundamental Theoretical Concepts}
%% -----------------------------------------------------------
The realization of a safe, AI-driven network administration agent relies on several interconnected disciplines. This section establishes each individually, building a progressive conceptual foundation.

\subsection{Evolution of Network Administration Paradigms}
Network infrastructure administration has historically evolved through three major paradigm shifts, driven by the growing complexity and scale of enterprise environments, as illustrated in Figure~\ref{fig:admin_evolution}.

\begin{figure}[htbp]
    \centering
    \begin{tikzpicture}[
        every node/.style={font=\sffamily\small},
        era/.style={rectangle, draw=black, thick, fill=white, text width=3.4cm, align=center, minimum height=2.6cm, rounded corners=5pt, inner sep=8pt},
        arrow/.style={->, very thick, draw=black!70, >=stealth},
        sublabel/.style={font=\sffamily\scriptsize, align=center, text width=3.2cm},
    ]
        \node[era] (manual) at (0, 0) {
            \textbf{Manual Administration}\\[5pt]
            \footnotesize CLI / GUI\\
            \footnotesize Direct commands\\
            \footnotesize Ad-hoc troubleshooting
        };
        \node[era] (iac) at (5.5, 0) {
            \textbf{Infrastructure-as-Code}\\[5pt]
            \footnotesize Ansible / Terraform\\
            \footnotesize Scripted playbooks\\
            \footnotesize Declarative / Imperative
        };
        \node[era] (ibn) at (11.0, 0) {
            \textbf{Intent-Based Networking}\\[5pt]
            \footnotesize Natural language intent\\
            \footnotesize AI-assisted translation\\
            \footnotesize Closed-loop verification
        };

        \draw[arrow] (manual.east) -- (iac.west);
        \draw[arrow] (iac.east) -- (ibn.west);

        \node[sublabel, below=0.4cm of manual] {High control,\\low scalability};
        \node[sublabel, below=0.4cm of iac] {Consistent, but\\rigid and static};
        \node[sublabel, below=0.4cm of ibn] {Adaptive and\\conversational};

    \end{tikzpicture}
    \caption{Evolution of Network Administration Paradigms}
    \label{fig:admin_evolution}
\end{figure}

\textbf{Manual Administration} relies on direct interaction via command-line interfaces (CLIs) or graphical interfaces. Although flexible for ad-hoc troubleshooting, manual processes fail to scale and place a severe cognitive burden on administrators, frequently resulting in configuration drift and costly human-induced outages~\cite{benson2009network,oppenheimer2003why}.

\textbf{Infrastructure-as-Code (IaC)} addressed these limitations by introducing a programmatic approach to resource provisioning. The IaC paradigm separates into two primary models: the \emph{imperative} model, where the administrator defines the exact sequence of commands required to reach the desired state (e.g., Ansible playbooks); and the \emph{declarative} model, where the administrator defines the desired final state, leaving the engine to compute the required delta~\cite{hashicorp_terraform} (e.g., Terraform). 

Despite these advancements, both IaC models remain static and cannot dynamically interpret the intent behind a conversational administrative request, nor adapt when live device state diverges from the template.

\textbf{Intent-Based Networking (IBN)} represents the emerging paradigm, where the administrator expresses high-level business or security intent in natural language, and the system autonomously translates, validates, and applies the corresponding configuration. 

According to Leivadeas and Falkner~\cite{leivadeas2022survey}, a mature IBN system operates on a closed-loop feedback model: intent is captured, translated into device configurations, the resulting state is verified against the intended outcome, and deviations are automatically remediated.

\subsection{Specific Challenges of Stateful Firewall Automation}
Firewall automation differs fundamentally from generic cloud resource orchestration. While provisioning a virtual machine is largely order-independent, managing security policies on a stateful firewall appliance introduces strict operational constraints that any automated agent must rigorously address~\cite{alshaer2004discovery}:

\begin{itemize}
    \item \textbf{Policy Ordering and Rule Shadowing:} Firewall Access Control Lists (ACLs) are evaluated in a top-down, first-match order. A broad permissive rule inserted at an incorrect index silently shadows specific restrictive policies, creating undetected security gaps~\cite{yuan2006fireman}.
    \item \textbf{Object Ambiguity and Entity Resolution:} Administrators refer to infrastructure by informal names (e.g., ``the mail server''). An automated agent must rigorously resolve these informal references to live, formal device objects — IP addresses, interface UUIDs, zone names — via active API queries, not static assumptions.
    \item \textbf{Configuration Drift and Out-of-Band Changes:} During security incidents, administrators apply emergency rules directly to the appliance CLI, immediately invalidating any static configuration template. Active state synchronization is therefore not optional — it is a fundamental architectural requirement.
    \item \textbf{Operational Criticality and Blast Radius:} A misconfigured virtual machine affects one host. A misconfigured firewall policy can disconnect entire subnets, sever corporate VPN gateways, or expose critical database segments to the public internet. This elevated blast radius demands rigid pre-execution verification gates before any change is applied.
\end{itemize}

This disconnect — between the abstract, business-driven language of the administrator and the rigid, syntax-heavy configuration commands required by the appliance — constitutes what is commonly referred to as the \textbf{Semantic Gap}. Bridging this gap without sacrificing execution predictability is the primary engineering challenge of modern network automation.

\subsection{Conversational AI: LLMs, RAG, and Safety Constraints}
The emergence of Large Language Models (LLMs) introduces the ability to declare administrative intent in natural language~\cite{boutaba2018comprehensive}. Rather than memorizing proprietary syntax, an administrator can simply describe what they need, and the model translates this into machine-readable instructions.

However, deploying a generic LLM as a direct execution agent over critical infrastructure introduces unacceptable risks. LLMs are inherently \emph{stochastic}: the same prompt may yield different outputs across invocations, and models are known to \emph{hallucinate} — generating plausible-sounding but factually incorrect values such as non-existent IP addresses, deprecated API fields, or invalid interface names~\cite{wang2024netconfeval}. Furthermore, a static language model has no awareness of the live state of the target device.

Two core technical safeguards are therefore mandatory to bridge the gap between the conversational power of LLMs and the rigid safety requirements of network security:

\begin{itemize}
    \item \textbf{Retrieval-Augmented Generation (RAG):} RAG grounds the language model's outputs by coupling it with a curated vector database \cite{lewis2020rag}. Before generating a response, the model retrieves the most semantically relevant technical documentation fragments and uses them as context. In network operations, this ensures that generated API payloads conform to documented schemas and that referenced objects exist on the live device~\cite{shajarian2024rag}.

    \item \textbf{Human-in-the-Loop (HITL) Validation:} Even a grounded LLM cannot guarantee the absence of policy-level errors in a live network context. The HITL paradigm~\cite{gomez2024hitl} positions the human administrator as the mandatory gatekeeper of all state mutations. A deterministic software layer intercepts the AI-generated payload, validates its structural and semantic correctness, presents a verifiable configuration diff to the operator, and awaits explicit approval before any command reaches the physical appliance.
\end{itemize}

Figure~\ref{fig:rag_pipeline} illustrates how RAG enrichment operates conceptually within the AI pipeline, grounding the model's generation with live context before a payload is produced.

\begin{figure}[htbp]
    \centering
    \begin{tikzpicture}[
        every node/.style={font=\sffamily\small},
        box/.style={rectangle, draw=black, thick, fill=white, text width=2.8cm, align=center, minimum height=1.1cm, rounded corners=4pt},
        dbbox/.style={cylinder, draw=black, thick, fill=gray!8, text width=2.2cm, align=center, minimum height=1.0cm, shape border rotate=90, aspect=0.25},
        arrow/.style={->, thick, draw=black, >=stealth},
        label/.style={font=\sffamily\scriptsize},
    ]

        \node[box] (query) at (0, 0) {\textbf{User Query}\\Natural language prompt};
        \node[dbbox] (vdb) at (5.5, -2.2) {Vector\\Database};
        \node[box] (context) at (5.5, 0) {\textbf{Retriever}\\Semantic search};
        \node[box] (llm) at (11.0, 0) {\textbf{LLM}\\Augmented\\generation};
        \node[box] (output) at (11.0, -2.5) {\textbf{Output}\\Grounded JSON\\payload};

        \draw[arrow] (query) -- node[above, label]{Query embedding} (context);
        \draw[arrow, transform canvas={xshift=0.3cm}] (context.south) -- node[right, label]{Top-k docs} (vdb.north);
        \draw[arrow, transform canvas={xshift=-0.3cm}] (vdb.north) -- node[left, label]{Context chunks} (context.south);
        \draw[arrow] (context) -- node[above, label]{Query + Context} (llm);
        \draw[arrow] (llm) -- (output);

    \end{tikzpicture}
    \caption{RAG Pipeline for LLM Grounding}
    \label{fig:rag_pipeline}
\end{figure}

%% -----------------------------------------------------------
\section{Comparative Analysis of Orchestration Models}
%% -----------------------------------------------------------
Having established the theoretical building blocks, it is necessary to contrast the proposed hybrid approach against existing orchestration paradigms before detailing its architecture. This comparison motivates the architectural direction of this project and establishes its unique contribution. Table~\ref{tab:orchestration_models} presents a structured comparison across four critical operational dimensions.

\begin{table}[htbp]
    \small
    \centering
    \renewcommand{\arraystretch}{1.4}
    \begin{tabular}{|p{2.8cm}|p{3.3cm}|p{3.8cm}|p{4.5cm}|}
        \hline
        \textbf{Feature} & \textbf{Static IaC (e.g., Terraform)} & \textbf{Pure LLM Agent} & \textbf{Hybrid Agent (Proposed)} \\
        \hline
        \textbf{Interface} & Code / CLI only & Conversational natural language & Conversational NL with live grounded context \\
        \hline
        \textbf{Safety Model} & Deterministic schema validation & Probabilistic, hallucination-prone~\cite{wang2024netconfeval} & Deterministic validation gate with mandatory HITL~\cite{gomez2024hitl} \\
        \hline
        \textbf{State Sync} & Static file, prone to drift & None --- blind execution & Active live API synchronization with pre-execution snapshots \\
        \hline
        \textbf{Recovery} & Manual destroy and re-apply & None (no rollback) & Automated transaction rollback on detected state drift (custom architecture) \\
        \hline
        \textbf{Adaptability} & None (rigid playbooks) & High (but unsafe) & High and safe (intent-routing with safety gates) \\
        \hline
    \end{tabular}
    \caption{Comparison of Network Orchestration Models}
    \label{tab:orchestration_models}
\end{table}

The hybrid model uniquely combines the conversational flexibility of LLM interfaces with the execution-time guarantees of a deterministic control plane. This synthesis motivates the architectural pattern formally defined in the following section.

%% -----------------------------------------------------------
\section{Theoretical Necessity of a Hybrid Architecture}
%% -----------------------------------------------------------
The preceding theoretical analysis converges on a single architectural necessity: the \textbf{Hybrid AI Network Orchestration Model}. This section formalizes the concept of a hybrid model, detailing why the isolation of generative inference from rule-based execution is critical for network safety.

\subsection{The Requirement for Execution Isolation}
The integration of generative models into critical infrastructure mandates the coexistence of two operationally distinct domains:

\begin{itemize}
    \item \textbf{Probabilistic Reasoning:} The capacity to parse unstructured natural language, retrieve organizational context, and formulate intended configurations. This domain relies on non-deterministic language models.
    \item \textbf{Deterministic Execution:} The capacity to validate syntactic correctness, enforce security policies, and apply state mutations to physical hardware. This domain must rely on rigid, verifiable logic.
\end{itemize}

The core theoretical necessity is that these two domains must never merge. They must be separated by a strict \textbf{Isolation Perimeter}: an architectural boundary that prevents any AI-generated artifact from reaching the network device without first undergoing programmatic verification and explicit human approval. This requirement dictates the transition from standalone AI chat interfaces to a hybrid, defensively engineered orchestration model.

\subsection{Redefining the Administrator's Role}
The adoption of a hybrid orchestration agent fundamentally redefines the operational responsibilities of the human administrator. Rather than eliminating the human from the process, the system strategically re-assigns tasks to maximize both operational velocity and safety guarantees. Table~\ref{tab:roles_actors} formally defines these reassigned responsibilities.

\begin{table}[htbp]
    \small
    \centering
    \renewcommand{\arraystretch}{1.5}
    \begin{tabular}{|p{4.0cm}|p{5.5cm}|p{5.0cm}|}
        \hline
        \textbf{Responsibility} & \textbf{Traditional Manual Admin} & \textbf{With AI Orchestration Agent} \\
        \hline
        \textbf{Intent Formulation} & Translates intent into CLI syntax manually & Expresses intent in natural language \\
        \hline
        \textbf{Configuration Drafting} & Writes commands directly & Delegated entirely to the AI Agent \\
        \hline
        \textbf{Safety Verification} & Performed manually (error-prone) & Automated by the Validation Engine \\
        \hline
        \textbf{Change Approval} & Implicit (execution = approval) & Explicit HITL confirmation required \\
        \hline
        \textbf{State Recovery} & Manual diagnosis and rollback & Automatic snapshot-based rollback \\
        \hline
        \textbf{Audit Trail} & Depends on manual logging discipline & Cryptographically appended audit log \\
        \hline
    \end{tabular}
    \caption{Operational Responsibilities Distribution}
    \label{tab:roles_actors}
\end{table}

The administrator's role is elevated from low-level syntax specialist to high-level intent decision-maker and safety authority. The AI agent absorbs the repetitive, syntactically intensive work, while the human retains ultimate authority over every state mutation applied to the live network.

%% -----------------------------------------------------------
\section{Comparative Analysis of Technical Solutions}
%% -----------------------------------------------------------
The choice of tools and technologies constitutes a critical engineering decision for each functional layer of the agent. This section analyzes the principal solutions available on the market and formally justifies the selection made for this project.

\subsection{Large Language Models (LLMs)}
The NLU engine requires a model capable of strictly following structured system instructions, producing syntactically correct JSON payloads, and operating with minimal latency for real-time conversational interactions. Three representative models were evaluated.

\begin{table}[htbp]
    \small
    \centering
    \renewcommand{\arraystretch}{1.4}
    \begin{tabular}{|p{3.0cm}|p{3.5cm}|p{3.5cm}|p{3.8cm}|}
        \hline
        \textbf{Criterion} & \textbf{OpenAI GPT-4} & \textbf{Meta Llama 3 (8B)} & \textbf{Mistral (mistral-small)} \\
        \hline
        \textbf{Deployment Mode} & Cloud API only & Local self-hosted & Cloud API or local \\
        \hline
        \textbf{JSON Schema Adherence} & Excellent (native function calling) & Requires extensive prompt engineering & Native strict JSON output mode \\
        \hline
        \textbf{API Response Latency} & Moderate (large model overhead) & Fast (GPU-dependent) & Very fast via lightweight API \\
        \hline
        \textbf{Cost Model} & High per-token pricing & Free (hardware cost only) & Highly cost-effective \\
        \hline
    \end{tabular}
    \caption{Comparison of LLMs for Intent Parsing}
    \label{tab:llm_comparison}
\end{table}

\textbf{Mistral AI (mistral-small-latest)} was selected. Its native JSON output mode enforces structured schema adherence at the model level, eliminating the risk of malformed payloads reaching the validation gate. Its very low inference latency makes it well-suited for real-time conversational interactions, and its cost-effective API pricing is appropriate for an academic engineering prototype requiring numerous validation test cycles.

\subsection{Vector Databases for the RAG Pipeline}
The RAG module requires a vector database capable of storing embedded documentation fragments and device object references, supporting fast semantic similarity search with no dependency on external cloud infrastructure.

\begin{table}[htbp]
    \small
    \centering
    \renewcommand{\arraystretch}{1.4}
    \begin{tabular}{|p{3.0cm}|p{3.5cm}|p{3.5cm}|p{3.8cm}|}
        \hline
        \textbf{Criterion} & \textbf{Pinecone} & \textbf{Qdrant} & \textbf{ChromaDB} \\
        \hline
        \textbf{Architecture} & Managed Cloud SaaS & Rust-based client/server & Python-native, embeddable \\
        \hline
        \textbf{Deployment} & External cloud, API keys required & Self-hosted via Docker & Zero-config embedded or containerized \\
        \hline
        \textbf{Data Sovereignty} & Data leaves the local network & Can be fully self-hosted & Fully local and persistent \\
        \hline
        \textbf{Python Integration} & REST API client library & REST API client library & Direct Python package import \\
        \hline
    \end{tabular}
    \caption{Comparison of Vector Databases}
    \label{tab:vector_db}
\end{table}

\textbf{ChromaDB} was selected for its Python-native architecture. It can operate in an embedded mode alongside the orchestrator or as a lightweight containerized service, requiring no external cloud infrastructure. All indexed FortiOS API documentation and device object references remain within the local environment, guaranteeing full data sovereignty for sensitive network inventory data.

\subsection{Backend Execution Framework}
The deterministic execution plane must dynamically construct, validate, and dispatch REST API calls to the FortiGate appliance based on the AI-generated payload received at runtime.

\begin{table}[htbp]
    \small
    \centering
    \renewcommand{\arraystretch}{1.4}
    \begin{tabular}{|p{3.0cm}|p{3.5cm}|p{3.5cm}|p{3.8cm}|}
        \hline
        \textbf{Criterion} & \textbf{Ansible (RedHat)} & \textbf{Nornir} & \textbf{Custom Python (httpx)} \\
        \hline
        \textbf{Execution Model} & Static YAML playbooks & Pluggable task execution & Synchronous/async HTTP client \\
        \hline
        \textbf{Dynamic JSON Handling} & Poor (static variable templating) & Moderate (requires plugin) & Excellent (full Python runtime) \\
        \hline
        \textbf{REST API Throughput} & Low (SSH-based overhead) & High (concurrent workers) & High (connection pooling, async I/O) \\
        \hline
        \textbf{Programmatic Flexibility} & Limited (no runtime branching logic) & Good & Full (native Python control flow) \\
        \hline
    \end{tabular}
    \caption{Comparison of Backend Frameworks}
    \label{tab:backend_frameworks}
\end{table}

A \textbf{custom Python orchestration engine using the httpx library} was selected as the REST client layer. While dedicated network automation frameworks like Ansible or Nornir excel at bulk configuration templating, they introduce unnecessary overhead when bridging dynamic AI JSON outputs directly to REST APIs. Python provides complete programmatic flexibility. The \texttt{httpx} library enables both synchronous and asynchronous HTTPS calls to the FortiOS REST API, supporting connection pooling and configurable timeouts — critical for reliably detecting management connectivity loss during post-execution verification.

\subsection{User Interface for the Security Operations Dashboard}
The Human-in-the-Loop confirmation workflow requires a dashboard interface that presents configuration diffs clearly, collects explicit operator approval, and provides real-time visibility into agent state.

\begin{table}[htbp]
    \small
    \centering
    \renewcommand{\arraystretch}{1.4}
    \begin{tabular}{|p{3.0cm}|p{3.5cm}|p{3.5cm}|p{3.8cm}|}
        \hline
        \textbf{Criterion} & \textbf{React / Node.js} & \textbf{Gradio} & \textbf{Streamlit} \\
        \hline
        \textbf{Implementation Language} & JavaScript / TypeScript & Python & Python \\
        \hline
        \textbf{Development Velocity} & Low (requires dedicated frontend) & Very high (limited customization) & High (flexible layout controls) \\
        \hline
        \textbf{Backend Integration} & Requires REST API bridge layer & Direct Python integration & Direct Python integration \\
        \hline
        \textbf{UI Customization} & Unlimited & Minimal & High (custom CSS, dynamic layouts) \\
        \hline
    \end{tabular}
    \caption{Comparison of UI Frameworks}
    \label{tab:ui_frameworks}
\end{table}

\textbf{Streamlit} was selected as the dashboard framework. It allows the rapid construction of interactive, Python-native web applications with direct access to the backend validation and session management logic. Its support for custom CSS injection, dynamic column layouts, and dark-mode theming produces a professional, SOC-appropriate dashboard aesthetic without requiring a separate JavaScript codebase or a REST API bridge layer.

\subsection{Target Network Appliance: The Firewall Ecosystem}
The target appliance must expose a comprehensive, programmatically accessible API permitting granular manipulation of security policies, address objects, NAT rules, and routing entries — without requiring a separate, mandatory management server.

\begin{table}[htbp]
    \small
    \centering
    \renewcommand{\arraystretch}{1.4}
    \begin{tabular}{|p{3.0cm}|p{3.5cm}|p{3.5cm}|p{3.8cm}|}
        \hline
        \textbf{Criterion} & \textbf{Cisco ASA} & \textbf{Palo Alto Networks} & \textbf{Fortinet FortiGate} \\
        \hline
        \textbf{API Accessibility} & Legacy REST API (often requires ASDM) & Excellent XML/REST API & Modern, granular REST API \\
        \hline
        \textbf{Management Model} & Standalone or FMC & Often relies on Panorama & Fully standalone capable \\
        \hline
        \textbf{Local Testing VM} & Available (ASAv) & Heavy resource requirements & Lightweight FortiOS VM available \\
        \hline
    \end{tabular}
    \caption{Comparison of Firewall Appliances}
    \label{tab:firewall_appliances}
\end{table}

\textbf{Fortinet FortiGate} running \textbf{FortiOS} was selected as the target appliance. While competitors like Palo Alto offer excellent APIs, they frequently rely on centralized management platforms (like Panorama) for complex orchestration. The FortiOS REST API provides a complete, well-documented interface for managing all network objects via token-authenticated HTTPS requests directly against the standalone appliance. This makes it ideally suited for an orchestration agent that must synchronize live state and apply precise mutations in real time. Furthermore, Fortinet's provision of a free FortiGate Virtual Machine edition enabled the establishment of a fully isolated local hypervisor environment for development and validation.

%% -----------------------------------------------------------
\section{Chapter Conclusion}
%% -----------------------------------------------------------
This chapter has progressively established the complete theoretical and technological foundation required for the realization of the hybrid AI orchestration agent. Starting from the cognitive burden of manual firewall administration and the rigid limitations of static IaC frameworks, the chapter built the case for a new operational paradigm grounded in Intent-Based Networking principles and AI-assisted natural language understanding.

The generic anatomy and functioning of a Hybrid AI Orchestration Agent were formally defined and illustrated: a system structured around four functional components operating across a strict Trust Boundary, through which every administrative request must travel from natural language input to deterministic, human-approved execution. The redistribution of operational responsibilities between administrator and agent was also formalized, highlighting the elevation of the human role from syntax specialist to intent authority and safety gatekeeper.

Finally, the comparative analysis of available technologies culminated in the formal justification of the selected stack: Mistral AI for intent parsing, ChromaDB for grounded RAG retrieval, Python with httpx for deterministic execution, Streamlit for the HITL dashboard, and FortiGate for the target appliance ecosystem.

With this foundation established, the following chapter will present the detailed system design and architectural blueprint of the implemented agent, translating these theoretical constructs into concrete software components, UML models, and validation mechanisms.
